\clearpage\section{MOS 電晶體電路圖}
下圖為可編輯的向量電路圖。同一電路的同名節點代表直接連接；獨立子電路單元的名稱為局部接腳，實際連接依階層圖。B 為基體端，0 為全域接地。各圖明列對應 netlist 與範圍，量測激勵源依文字說明接入。
實際單元與階層對應 \texttt{01\_mux\_truth\_table.sp} 及 \texttt{02\_mux\_delay\_load.sp}.
\par
\begin{figure}[!ht]
\centering
\begin{circuitikz}[font=\small]
\node[pmos,bulk,rotate=-90] (TP) at (0,1) {};
\node[nmos,bulk,rotate=90] (TN) at (0,-1) {};
\draw (TP.D) -- (TP.D -| -2,0) -- (TN.D -| -2,0) -- (TN.D);
\draw (TP.S) -- (TP.S -| 2,0) -- (TN.S -| 2,0) -- (TN.S);
\draw (-2,0) node[circ]{} -- (-2.6,0) node[left]{a (D)};
\draw (2,0) node[circ]{} -- (2.6,0) node[right]{z (S)};
\draw (TP.G) -- ++(0,.5) node[above]{eb};
\draw (TN.G) -- ++(0,-.5) node[below]{en};
\draw (TP.bulk) -- ++(0,-.35) -- ++(.8,0) node[right,font=\scriptsize]{B=vp};
\draw (TN.bulk) -- ++(0,.35) -- ++(.8,0) node[right,font=\scriptsize]{B=vn};
\node[font=\scriptsize] at (-.8,1.4){TP};
\node[font=\scriptsize] at (-.8,-1.4){TN};
\end{circuitikz}
\caption{兩種傳輸閘架構共用的 MOS 傳輸單元。}
\end{figure}
\par
\par
\begin{figure}[!ht]
\centering
\begin{circuitikz}[font=\small]
\node[pmos,bulk] (MP) at (0,2.5) {};
\node[anchor=west,font=\scriptsize] at ($(MP.center)+(.55,.40)$) {MP};
\draw (MP.bulk) -- ++(.30,0) node[right,font=\scriptsize] {B=vp};
\node[nmos,bulk] (MN) at (0,0.2) {};
\node[anchor=west,font=\scriptsize] at ($(MN.center)+(.55,.40)$) {MN};
\draw (MN.bulk) -- ++(.30,0) node[right,font=\scriptsize] {B=vn};
\draw (MP.G) -- (-1,2.5) -- (-1,.2) -- (MN.G); \draw (-1,1.35) node[circ]{} -- (-1.6,1.35) node[left]{a}; \draw (MP.D |- 0,1.35) node[circ]{} -- (2,1.35) node[right]{z};
\draw (MP.S) -- ++(0,0.5) node[above,font=\small] {\texttt{vp}};
\draw (MP.D) -- (MN.D) ;
\draw (MN.S) -- ++(0,-0.5) node[below,font=\small] {\texttt{vn}};
\end{circuitikz}
\caption{輸入／輸出與局部選擇訊號反相器。}
\end{figure}
\par
inv 單元的 $W_P=3.6k\,\mu$m、$W_N=1.8k\,\mu$m、$L=0.25\,\mu$m；tg 使用相同尺寸。tg 的兩顆汲極均為 a、源極均為 z，但傳輸路徑為雙向。圖中 TP／TN 對應 tg 子電路內的 MP／MN。PMOS 基體固定接 vp，NMOS 基體固定接 vn，並非接訊號端。NMOS 閘極為 en，PMOS 閘極為 eb；同名節點將兩傳輸電晶體並聯。
\clearpage
\par
\begin{figure}[!ht]
\centering
\begin{circuitikz}[font=\small]
\node[draw,rounded corners=2pt,align=center,minimum width=1.6cm,minimum height=.9cm] (T0) at (1,3) {XFT0\\TG: fb0};
\draw (-1,3) node[left]{f\_in0} -- (T0.west);
\node[draw,rounded corners=2pt,align=center,minimum width=1.6cm,minimum height=.9cm] (T1) at (1,1) {XFT1\\TG: s0};
\draw (-1,1) node[left]{f\_in1} -- (T1.west);
\node[draw,rounded corners=2pt,align=center,minimum width=1.6cm,minimum height=.9cm] (T2) at (1,-1) {XFT2\\TG: fb0};
\draw (-1,-1) node[left]{f\_in2} -- (T2.west);
\node[draw,rounded corners=2pt,align=center,minimum width=1.6cm,minimum height=.9cm] (T3) at (1,-3) {XFT3\\TG: s0};
\draw (-1,-3) node[left]{f\_in3} -- (T3.west);
\node[draw,rounded corners=2pt,align=center,minimum width=1.6cm,minimum height=.9cm] (T4) at (5,2) {XFT4\\TG: fb1};
\node[draw,rounded corners=2pt,align=center,minimum width=1.6cm,minimum height=.9cm] (T5) at (5,-2) {XFT5\\TG: s1};
\draw (T0.east) -- (3,3) -- (3,2) node[circ,label=above right:fn0]{} -- (T4.west); \draw (T1.east) -| (3,2); \draw (T2.east) -- (3,-1) -- (3,-2) node[circ,label=above right:fn1]{} -- (T5.west); \draw (T3.east) -| (3,-2); \draw (T4.east) -- (7,2) -- (7,0) node[circ]{} -- (8,0) node[right]{of}; \draw (T5.east) -| (7,0);
\end{circuitikz}
\caption{最佳化 TG 樹：TG 後的標示為 en，eb 為其局部反相；六個單元皆 k=8。}
\end{figure}
\par
fb0／fb1 由兩個 inv 對 s0／s1 反相產生，使用 vf／0。此樹的所有 TG 基體皆接 vf／0；CLf 由 of 接 0。沒有輸出反相器：六個 TG 加兩個選擇反相器共 16 顆 MOS。零伏輸入探棒 VDfi 的正端接 f\_ini、負端接 ai；g／c 架構的探棒採相同接法。
\clearpage
\par
\begin{figure}[!ht]
\centering
\begin{circuitikz}[font=\small]
\node[draw,rounded corners=2pt,align=center,minimum width=1.8cm,minimum height=.9cm] (A) at (0,0) {g\_ini\\$i=0\ldots3$};
\node[draw,rounded corners=2pt,align=center,minimum width=1.8cm,minimum height=.9cm] (B) at (3,0) {XGIi: inv};
\node[draw,rounded corners=2pt,align=center,minimum width=1.8cm,minimum height=.9cm] (C) at (6,0) {XGTia\\TG: $S_{1i}$};
\node[draw,rounded corners=2pt,align=center,minimum width=1.8cm,minimum height=.9cm] (D) at (9,0) {XGTib\\TG: $S_{0i}$};
\draw[-{Latex}] (A.east) -- (B.west);
\draw[-{Latex}] (B.east) -- (C.west) node[midway,above]{gai};
\draw[-{Latex}] (C.east) -- (D.west) node[midway,above]{gni};
\draw (D.east) -- (11,0) node[right]{gz};
\node[draw,rounded corners=2pt,align=center,minimum width=1.2cm,minimum height=.9cm] (E) at (3,-2) {gz};
\node[draw,rounded corners=2pt,align=center,minimum width=1.8cm,minimum height=.9cm] (F) at (6,-2) {XGO: inv};
\node[draw,rounded corners=2pt,align=center,minimum width=1.2cm,minimum height=.9cm] (G) at (9,-2) {og};
\draw[-{Latex}] (E.east) -- (F.west);
\draw[-{Latex}] (F.east) -- (G.west);
\end{circuitikz}
\caption{講義 TG 架構：四條索引支路共接 gz，再由單一輸出反相器產生 og。}
\end{figure}
\par
當 $i\ge2$ 時 $S_{1i}=s1$，否則為 gb1；奇數 i 時 $S_{0i}=s0$，否則為 gb0。eb 使用相反的選擇文字。所有單元 k=1，電源為 vg／0，另含兩個局部選擇反相器。四個分支輸入與共同輸出皆有反相器，含選擇反相共 30 顆 MOS。
\par
\begin{figure}[!ht]
\centering
\begin{circuitikz}[font=\small]
\node[draw,rounded corners=2pt,align=center,minimum width=3.5cm,minimum height=.9cm] (N) at (0,0) {$(S_{1i},S_{0i},c\_ini)$};
\node[draw,rounded corners=2pt,align=center,minimum width=2.2cm,minimum height=.9cm] (P) at (4,0) {XCN$i$: NAND3\\$i=0\ldots3$};
\node[draw,rounded corners=2pt,align=center,minimum width=3.5cm,minimum height=.9cm] (Q) at (8,0) {XCO: NAND4\\$(cn0,cn1,cn2,cn3)$};
\draw[-{Latex}] (N.east) -- (P.west);
\draw[-{Latex}] (P.east) -- (Q.west) node[midway,above]{cni};
\draw (Q.east) -- (10.5,0) node[right]{oc};
\end{circuitikz}
\caption{靜態 CMOS MUX 階層：四個 NAND3 輸出進入一個 NAND4。}
\end{figure}
\par
此處的反相選擇訊號為 cb0／cb1，由 vc／0 的局部反相器產生；靜態邏輯閘皆 k=1。og／oc／of 各有相同 cl 負載接 0。
\clearpage
\par
\begin{figure}[!ht]
\centering
\begin{circuitikz}[font=\small]
\node[pmos,bulk] (MP0) at (0.0,3) {};
\node[anchor=west,font=\scriptsize] at ($(MP0.center)+(.55,.40)$) {MP0};
\draw (MP0.bulk) -- ++(.30,0) node[right,font=\scriptsize] {B=vp};
\draw (MP0.G) -- ++(-.45,0) node[left,font=\small] {\texttt{a0}};
\draw (MP0.S) -- (MP0.S |- 0,4.2); \draw (MP0.D) -- (MP0.D |- 0,1.7);
\node[nmos,bulk] (MN0) at (3.8,0.0) {};
\node[anchor=west,font=\scriptsize] at ($(MN0.center)+(.55,.40)$) {MN0};
\draw (MN0.bulk) -- ++(.30,0) node[right,font=\scriptsize] {B=vn};
\draw (MN0.G) -- ++(-.45,0) node[left,font=\small] {\texttt{a0}};
\node[pmos,bulk] (MP1) at (3.8,3) {};
\node[anchor=west,font=\scriptsize] at ($(MP1.center)+(.55,.40)$) {MP1};
\draw (MP1.bulk) -- ++(.30,0) node[right,font=\scriptsize] {B=vp};
\draw (MP1.G) -- ++(-.45,0) node[left,font=\small] {\texttt{a1}};
\draw (MP1.S) -- (MP1.S |- 0,4.2); \draw (MP1.D) -- (MP1.D |- 0,1.7);
\node[nmos,bulk] (MN1) at (3.8,-2.2) {};
\node[anchor=west,font=\scriptsize] at ($(MN1.center)+(.55,.40)$) {MN1};
\draw (MN1.bulk) -- ++(.30,0) node[right,font=\scriptsize] {B=vn};
\draw (MN1.G) -- ++(-.45,0) node[left,font=\small] {\texttt{a1}};
\draw (MN0.S) -- (MN1.D) node[midway,above right,font=\small] {\texttt{n1}};
\node[pmos,bulk] (MP2) at (7.6,3) {};
\node[anchor=west,font=\scriptsize] at ($(MP2.center)+(.55,.40)$) {MP2};
\draw (MP2.bulk) -- ++(.30,0) node[right,font=\scriptsize] {B=vp};
\draw (MP2.G) -- ++(-.45,0) node[left,font=\small] {\texttt{a2}};
\draw (MP2.S) -- (MP2.S |- 0,4.2); \draw (MP2.D) -- (MP2.D |- 0,1.7);
\node[nmos,bulk] (MN2) at (3.8,-4.4) {};
\node[anchor=west,font=\scriptsize] at ($(MN2.center)+(.55,.40)$) {MN2};
\draw (MN2.bulk) -- ++(.30,0) node[right,font=\scriptsize] {B=vn};
\draw (MN2.G) -- ++(-.45,0) node[left,font=\small] {\texttt{a2}};
\draw (MN1.S) -- (MN2.D) node[midway,above right,font=\small] {\texttt{n2}};
\draw (MP0.S |- 0,4.2) -- (MP2.S |- 0,4.2) node[midway,above]{vp}; \draw (MP0.D |- 0,1.7) -- (MP2.D |- 0,1.7); \draw (MN0.D) -- (MN0.D |- 0,1.7) node[circ,label=above right:z]{};
\draw (MN2.S) -- ++(0,-0.4) node[below,font=\small] {\texttt{vn}};
\end{circuitikz}
\caption{實際 NAND3：三條 PMOS 並聯支路與透過 n1／n2 串聯的 NMOS。}
\end{figure}
\par
nand3 的 $W_P=3.6k\,\mu$m、$W_N=5.4k\,\mu$m、$L=0.25\,\mu$m。nand4 以相同 PMOS 並聯／NMOS 串聯方式擴充至四輸入，$W_N=7.2k\,\mu$m，PMOS 寬度與長度不變。所有 PMOS 基體／源極接 vp；NMOS 基體皆接 vn，但中間源極為串疊節點。因此四個 NAND3（24 顆）、一個 NAND4（8 顆）、兩個選擇反相器（4 顆），總共 36 顆 MOS。
